\documentclass[12pt,a4paper]{article}
\usepackage[T2A]{fontenc}
\usepackage[utf8]{inputenc}
\usepackage[russian]{babel}
\usepackage{amsmath}
\usepackage{geometry}
\usepackage{fancyvrb}
\usepackage{graphicx}
\geometry{left=2cm,right=2cm,top=2cm,bottom=2cm}
\begin{document}
\pagestyle{empty}
\begin{center}
    \setstretch{1.1}
    {\fontsize{12pt}{14pt}\selectfont
        МИНИСТЕРСТВО НАУКИ И ВЫСШЕГО ОБРАЗОВАНИЯ\\
        РОССИЙСКОЙ ФЕДЕРАЦИИ\\
        Федеральное государственное бюджетное образовательное учреждение\\
        высшего образования\\
        \textbf{АДЫГЕЙСКИЙ ГОСУДАРСТВЕННЫЙ УНИВЕРСИТЕТ}\\
    }
    \vspace{0.3cm}
    {\fontsize{12pt}{14pt}\selectfont
        Инженерно-физический факультет\\
        Кафедра автоматизированных систем обработки информации и\\
        управления\\
    }
\end{center}
\vspace{4.5cm}
\begin{center}
    \setstretch{1.3}
    {\fontsize{14pt}{16pt}\selectfont ОТЧЕТ ПО ПРАКТИКЕ\\}
    \vspace{0.5cm}
    {\fontsize{16pt}{18pt}\selectfont \textbf{Программная реализация численного метода}\\}
    {\fontsize{14pt}{16pt}\selectfont \textit{Текст из задания по варианту 8}\\}
    \vspace{0.8cm}
    {\fontsize{12pt}{14pt}\selectfont 1 курс, группа УТС\\}
\end{center}
\vspace{3cm}
\hfill
\begin{minipage}{0.48\textwidth}
    \setstretch{1.2}
    {\normalsize Выполнил:\\
    \underline{\hspace{3.5cm}} В. Д. Платоненко\\
    «\underline{\hspace{0.8cm}}» \underline{\hspace{2.2cm}} 2026 г.\\
    \vspace{0.6cm}
    Руководитель:\\
    \underline{\hspace{3.5cm}} С. В. Теплоухов\\
    «\underline{\hspace{0.8cm}}» \underline{\hspace{2.2cm}} 2026 г.}
\end{minipage}
\vfill 
\begin{center}
    {\normalsize Майкоп, 2026 г.}
\end{center}
\newpage
\section{Задание}
Реализовать алгоритм вычисления первой производной функции в точке тремя методами:
\begin{enumerate}
    \item Метод правой разности.
    \item Метод центральной разности.
    \item Метод левой разности.
\end{enumerate}
В качестве тестируемой функции использовать $f(x) = x^2$. Количество точек и шаг задаются пользователем.
\section{Теоретическая часть}
Для численного нахождения производной функции $f(x)$ в точке $x$ с шагом $dX$ применяются следующие формулы:
\begin{itemize}
    \item \textbf{Правая разностная производная (Метод 1):}
    \begin{equation}
        f'(x) \approx \frac{f(x + dX) - f(x)}{dX}
    \end{equation}
    \item \textbf{Центральная разностная производная (Метод 2):}
    \begin{equation}
        f'(x) \approx \frac{f(x + dX) - f(x - dX)}{2 \cdot dX}
    \end{equation}
    \item \textbf{Левая разностная производная (Метод 3):}
    \begin{equation}
        f'(x) \approx \frac{f(x) - f(x - dX)}{dX}
    \end{equation}
\end{itemize}
\section{Реализация и контроль ввода}
Программа написана на языке C++. Для обеспечения надежности реализован контроль вводимых данных:
\begin{itemize}
    \item Проверка типа данных при вводе $x$ и $dX$ (защита от ввода букв с помощью очистки буфера \texttt{cin.clear()} и \texttt{cin.ignore()}).
    \item Проверка значения шага $dX$ (шаг не должен быть равен нулю для избежания деления на ноль).
\end{itemize}
\newpage
\section{Исходный код программы}
Ниже представлен полный исходный код разработанной программы:
\begin{Verbatim}[frame=single, numbers=left, tabsize=4]
#include <iostream>
#include <cmath>
#include <clocale>
using namespace std;

double f(double x) {
    return x * x;
}

int main() {
    setlocale(LC_ALL, "Russian");
    double x, dX;
    
    while (true) {
        cout « "Введите число X = ";
        if (cin » x) {
            break;
        }
        else {
            cout « "Введено не число" « endl;
            cin.clear();
            cin.ignore(1000, '\n');
        }
    }
    while (true) {
        cout « "Введите ненулевое число dX = ";
        if (cin » dX && dX != 0) {
            break;
        }
        else {
            cout « "dX должно быть числом, но не 0" « endl;
            cin.clear();
            cin.ignore(1000, '\n');
        }
    }
    double Pravaya_formula = (f(x + dX) - f(x)) / dX;
    double Centralnaya_formula = (f(x + dX) - f(x - dX)) / (2 * dX);
    double Levaya_formula = (f(x) - f(x - dX)) / dX;
    
    cout « "Правая разностная производная = " « Pravaya_formula « endl;
    cout « "Центральная разностная производная = " « Centralnaya_formula « endl;
    cout « "Левая разностная производная = " « Levaya_formula « endl;
    
    return 0;
}
\end{Verbatim}
\newpage
\section{Результат работы программы}
\begin{figure}[h!]
\centering
\includegraphics[widht=0.8\textwidth]{изображение_2026-06-11_172754739.png}
\caption{Скриншот выполнения программы}
\label{fig:screenshot}
\end{figure}
\end{document}
